\documentclass[11pt]{article}
\usepackage[margin=1in]{geometry}
\usepackage{amsmath,amssymb,amsthm,mathtools,enumitem,booktabs}
\usepackage[colorlinks=true,linkcolor=blue,citecolor=blue,urlcolor=blue]{hyperref}

\newtheorem{theorem}{Theorem}\newtheorem{lemma}[theorem]{Lemma}
\newtheorem{proposition}[theorem]{Proposition}\newtheorem{conjecture}[theorem]{Conjecture}
\newtheorem{corollary}[theorem]{Corollary}
\theoremstyle{definition}\newtheorem{remark}[theorem]{Remark}
\newtheorem{finding}[theorem]{Finding}

\DeclareMathOperator{\ord}{ord}\DeclareMathOperator{\den}{den}\DeclareMathOperator{\lcm}{lcm}
\newcommand{\Z}{\mathbb{Z}}\newcommand{\Q}{\mathbb{Q}}\newcommand{\Fp}{\mathbb{F}_p}
\newcommand{\wt}{\widetilde}\newcommand{\dn}{d_n}
\newcommand{\zt}[1]{\zeta(#1)}

\title{The totally symmetric $\zeta(7)$ cellular family:\\
a verified recurrence, two period ladders, the value of $P_3$,\\
and a Dwork-congruence reduction of the central-binomial law}
\author{}
\date{18 July 2026}

\begin{document}
\maketitle

\begin{abstract}
We report the results of an intensive computational--theoretical campaign on
Brown--Zudilin cellular integrals at weights five and seven. For the totally
symmetric $\zeta(7)$ family on $\overline{\mathcal M}_{0,10}$ we reconstruct an
exact candidate recurrence for the leading coefficients $q_n$ (order $4$, degree
$19$, from $105$ computed terms; verified against all $74$ exact values), whose
palindromic characteristic polynomial
$\lambda^4-6340\lambda^3+67974\lambda^2-6340\lambda+1$ factors over
$\Q(\sqrt3)$ into two reciprocal quadratic sectors. The family splits into two
period ladders: conditionally, the weight-5 descent $(q_n,s_n,\hat P_n)$ rides
the recurrence, while the primitive $(P_n,I'_n)$ does not. Using the
filtration identity $I_n=I'_n+\zt2\,I''_n$ and the resulting conditional
values $s_3,\hat P_3$, we determine, conditional also on the denominator grid,
$P_3=23478462179525/(2^5 3^7)$ by a denominator-snap argument that requires
\emph{no} integral evaluation, verified independently ($99.75\%$ series
climb; rigorous tail bound). Under those hypotheses,
$\den(P_3)\mid d_3^7$: the decisive
$3$-adic cell shows no denominator anomaly. Conditional sector measurements
place both projections on the sector-B rate $0.0939$; with the stated
denominator cost, the standard criterion therefore cannot prove ``one of
$\zt5,\zt7$ is irrational''. The $\zt2$-elimination costs a factor
$\approx594$ in decay, the
archimedean mirror of the Betti lattice cost. On the weight-5 arithmetic side,
the primes $p\le n$ of the central-binomial divisibility
$\binom{2n}{n}\mid6A_n$ (Phase 2) are reduced to a single new congruence: a
vector-valued \emph{Dwork descent} $p^5(p_n/q_n)\equiv p_a/q_a \bmod p^{2-\kappa}$,
$a=\lfloor n/p\rfloor$, whose constants are the $n=1$ ladder ratios
$29/28=p_1/q_1$, $101/84=\wt p_1/q_1$; subject also to the required
integer-lattice normalization at exceptional primes, the general inequality
$\ord_p(p_n)\ge\kappa-5L$ is verified with zero failures. Every claim below is
labeled proven, verified, or conjectural.
\end{abstract}

\section{Setting and summary of status}

Brown--Zudilin \cite{BZ} construct cellular integrals $I_n$ on
$\overline{\mathcal M}_{0,N}$; at $N=10$, the totally symmetric convergent cell
$\sigma=(10,2,4,1,6,3,8,5,9,7)$ yields a rank-4 motive with weight-graded
pieces $\Q(0),\Q(-2),\Q(-5),\Q(-7)$ and linear forms
\[
I'_n=\tfrac{75}{4}q_n\zt7-3s_n\zt5-P_n,\qquad
I''_n=-9q_n\zt5+2s_n\zt3-\hat P_n,\qquad
I_n=I'_n+\zt2\,I''_n,
\]
with printed anchors $q_n=1,61,52921$; $s_n=0,300,261153$; $P_n=0,220,6021219/32$;
$\hat P_n=0,152,535857/4$ for $n=0,1,2$ \cite{BZ}. The leading coefficients
$q_n$ are computable for all $n$ by the McCarthy--Osburn--Straub diagonal
construction \cite{MOS}, verified against the anchors and extended here to
$n\le104$.

\smallskip\noindent\textbf{Status ledger.} Throughout, \emph{proven} = complete
proof; \emph{verified} = exact finite computation, no proof for all $n$;
\emph{conditional} = exact consequence of named hypotheses.

\begin{center}\small
\begin{tabular}{lll}
\toprule
Result & Status & Where \\
\midrule
$(W)$: $w_n\equiv\wt w_n\equiv0\ (p),\ n<p\le2n$ & \textbf{proven} (all $n$) & \cite{CBcert} \\
Phase 1 of $\binom{2n}{n}\mid6A_n$ & \textbf{proven} (all $n$) & \cite{CBcert} \\
$q_n$ recurrence (order 4, deg 19) & verified ($74$ terms) & \S\ref{sec:rec} \\
Descent membership of $s_n,\hat P_n$ & conditional/compatible with anchors & \S\ref{sec:ladders} \\
$s_3,\hat P_3$ exact values & conditional on membership & \S\ref{sec:ladders} \\
$P_3$ exact value & conditional (2 hypotheses) & \S\ref{sec:snap} \\
Sector measurement (decay rates) & conditional/measured & \S\ref{sec:sector} \\
Phase-2 assembly from (DWORK$_\Z$) & proven conditionally; verified data & \S\ref{sec:dwork} \\
(DWORK), incl. exceptional normalization & \textbf{open} (new) & \S\ref{sec:dwork} \\
\bottomrule
\end{tabular}
\end{center}

\section{The exact recurrence and its two sectors}\label{sec:rec}

\begin{finding}[verified candidate]\label{find:rec}
The computed sequence $q_n$ of the totally symmetric $M_{0,10}$ cell is
annihilated, throughout the available exact range, by a linear recurrence of
order $4$ with polynomial coefficients of degree $19$
(leading/trailing coefficients palindromic, $7381728\cdot[1,-6340,67974,-6340,1]$),
reconstructed from $105$ terms by $3$-prime CRT ($\approx2^{125}$) with exact
rational lifting, and certified to annihilate all $74$ exact values of $q_n$
(all $70$ testable offsets vanish identically; independently re-verified).
The characteristic polynomial
\[
\chi(\lambda)=\lambda^4-6340\lambda^3+67974\lambda^2-6340\lambda+1
=\bigl(\lambda^2-a\lambda+1\bigr)\bigl(\lambda^2-b\lambda+1\bigr),
\qquad a,b=3170\pm1824\sqrt3,
\]
factors over $\Q(\sqrt3)$ (check: $a+b=6340$, $2+ab=67974$). The four roots
form two reciprocal pairs (Poincar\'e-duality pairing):
\emph{sector A} $\{6329.2605\ldots,\ 1.57996\ldots\times10^{-4}\}$ and
\emph{sector B} $\{10.6454\ldots,\ 0.0939374\ldots\}$.
\end{finding}

The order matches the motivic rank ($\zt3$: order 2; $\zt5$: order 3; $\zt7$:
order 4), while the degree jump ($9\to19$ from weight 5) shows arithmetic
complexity growing faster than rank. Annihilation of finitely many terms is
not a proof that the operator governs $q_n$ for all $n$; a creative-telescoping
certificate was attempted on two representations (an $8$-variable low-coupling
window model and the $4$-fold hypergeometric-sum form of \S\ref{sec:snap}) and
in both cases hit the geometric elimination wall (successive eliminations
$\ge3$h each); checkpoints are preserved. The recurrence's status therefore
remains \emph{verified}, with overwhelming structural support (motivic order
match, palindromy, two independent window representations agreeing on all
computed terms).

\section{Two period ladders}\label{sec:ladders}

\begin{finding}[conditional, compatible with all available anchors]\label{find:ladders}
The weight-5 descent data $(q_n,s_n,\hat P_n)$ --- equivalently the form
$I''_n$ --- ride the order-4 operator: the three known values of each,
augmented by the operator's trailing-coefficient degeneracy at the initial
offset (which also re-derives $q_3=94357501$ from $q_0,q_1,q_2$, an internal
consistency check), determine
\[
s_3=\frac{1396906795}{3},\qquad
\hat P_3=\frac{232175579999}{972}\quad(972=2^2 3^5).
\]
The primitive data $(P_n, I'_n)$ do \emph{not} satisfy the operator:
propagation is inconsistent (the naive $P_3$ candidate even acquires a
spurious prime factor $107$). Weight seven thus splits into two holonomic
species --- the descended and the primitive ladder --- which weight five does
not separate.
\end{finding}

The values $s_3,\hat P_3$ are exact \emph{conditional on descent membership}
(that the actual $s_n,\hat P_n$ satisfy the operator for all $n$); membership
is compatible with every available anchor and is suggested by the motivic rank
structure, but a period-level telescoping certificate is the outstanding
closure (see \S\ref{sec:outlook}).

\section{The value of $P_3$ by denominator snap}\label{sec:snap}

The full cellular integral admits the exact all-positive $4$-fold series
representation of \cite{BarnesNotes} (verified against the anchors), hence
$I_n>0$ and $I_{n+1}<I_n$ rigorously. The observed values collapse:
$I_0=3.5554\ldots$, $I_1=3.2071\ldots\times10^{-5}$,
$I_2=1.05313\ldots\times10^{-9}$ --- the full form is a near-perfect
cancellation $I'\approx-\zt2 I''$ (four orders at $n=2$).

\begin{proposition}[snap; conditional on (i)--(ii)]\label{prop:snap}
Assume \textup{(i)} $\den(P_3)\mid 12\,d_3^7$ and \textup{(ii)} descent
membership (so $s_3,\hat P_3$ are as above). Then
\[
P_3=\frac{23478462179525}{69984},\qquad
\den(P_3)=69984=2^5\,3^7 .
\]
\end{proposition}

\begin{proof}[Proof (no direct evaluation of $I_3$)]
By the filtration identity, $P_3=\bigl[\tfrac{75}{4}q_3\zt7-3s_3\zt5
+\zt2 I''_3\bigr]-I_3$, where the bracket is computable to arbitrary precision
from (ii) and $I''_3=-9q_3\zt5+2s_3\zt3-\hat P_3$. The grid $(1/12d_3^7)\Z$ of
hypothesis (i) has spacing $h=2.977\times10^{-7}$. Certified rational
enclosures for the zeta values place the bracket strictly between the stated
grid point and that grid point plus $h/2$. Independently, positivity together
with the rigorous series majorant gives
$0<I_3\le3.09\times10^{-9}<h/2$. Hence $P_3$ is the unique grid point in the
permitted one-sided interval below the bracket, and it is the stated rational.
The decimal proximity is not the certification: it only indicates which
short certified enclosure to check.
\end{proof}

\noindent\textbf{Independent verification.} The implied value
$I_3=5.6299224184893\times10^{-14}$ was tested against the all-positive
series: partial sums (rigorous lower bounds) climb to $99.75\%$ of it without
ever exceeding it, and power-law extrapolation lands within $0.04\%$; a proved
separable majorant gives the rigorous bound $I_3\le3.09\times10^{-9}<h/2$,
making the half-window premise unconditional (the snap still assumes
(i)--(ii)). Consistency locks: the
residual is positive as required; $I_3/I_2=5.35\times10^{-5}$ continues the
ratio sequence $9.0\times10^{-6},\,3.28\times10^{-5}$ monotonically toward the
sector-A root $1.58\times10^{-4}$, as duality predicts for the full form; a
wrong grid hypothesis would land within $5.6\times10^{-14}$ of a grid point;
under a uniform-phase heuristic this has probability $\sim2\times10^{-7}$.

\begin{corollary}[conditional]\label{cor:noanomaly}
$\den(P_3)=2^53^7$ divides $d_3^7=2^73^7$: the naive inclusion
$d_3^7P_3\in\Z$ holds at $n=3$, with $\ord_3$ tight and $\ord_2$ slack $2$. At
its first $3$-adic opportunity the weight-7 primitive ladder shows \emph{no}
$12$-type correction --- consistent with the $n\le2$ data and with
$\den(\hat P_3)=2^23^5$ on the descent ladder: the anomalous correction factor
measured at weight 5 \cite{Denoms} does not recur in this family through $n=3$.
\end{corollary}

\section{Sector measurement and the Diophantine verdict}\label{sec:sector}

For a weight-7 form with denominator growth $d_n^{7+o(1)}$, irrationality-type
conclusions require decay below $e^{-7}=9.119\times10^{-4}$. The two sub-unit
roots straddle this threshold: sector A's $1.58\times10^{-4}$ passes (margin
$e^{1.753}$), sector B's $0.0939$ fails by two orders.

\begin{finding}[measured, conditional on the recurrence and descent membership]\label{find:sector}
Propagating $I''_n$ through the exact operator at $300$-digit precision from
its four exact initial values gives ratios increasing monotonically
$0.0042,\ 0.0181,\ 0.0306,\ \ldots,\ 0.0833\ (n=29)$ toward $0.0939$: the
\emph{descended} ladder rides sector B. The filtration then locks the
primitive projection: since $I_n$ (sector A, ratios
$9.0\times10^{-6},3.3\times10^{-5},5.4\times10^{-5}\to1.58\times10^{-4}$) is
negligible against $\zt2 I''_n$, we get $I'_n\approx-\zt2I''_n$, and indeed
$I'_3/I'_2=0.0306408$ against $I''_3/I''_2=0.0306408$ (difference
$2.8\times10^{-8}$). \emph{Both projections ride sector B.}
\end{finding}

\begin{corollary}[conditional on the recurrence, sector assignment, descent
membership, and stated denominator cost]
The standard simultaneous-linear-form criterion applied to these two
projections cannot prove ``at least one of $\zt5,\zt7$ is irrational'' under a
$d_n^7$-type denominator cost: the
required decay $e^{-7}$ is missed by the factor $0.0939/e^{-7}\approx103$.
The $\zt2$-elimination costs a decay factor $0.0939/1.58\times10^{-4}
\approx594$ --- an archimedean analogue of the Betti elimination cost of
\cite{Denoms}. The surviving Diophantine route is asymmetric: optimize
$\mathcal M(a)=-\log\rho_{\mathrm{proj}}(a)-\kappa(a)$ over the parameter
orbit (decay and denominator cost move together; positive $\mathcal M$ is the
Diophantine region).
\end{corollary}

\section{Phase 2 of the central-binomial law: reduction to a Dwork
congruence}\label{sec:dwork}

We now return to weight 5. With Zudilin's sequences $q_n,p_n,\wt p_n$
\cite{Zud02} and $A_n=2d_n^5p_n\in\Z$, the corrected symmetric law
$12d_n^5P_n\in\Z$ is equivalent to $\binom{2n}{n}\mid6A_n$, i.e.
$\ord_p(p_n)\ge\kappa-5L$ at every prime $p\ge5$, where
$\kappa=\ord_p\binom{2n}{n}$ (Kummer's carry count) and
$L=\ord_p d_n=\lfloor\log_p n\rfloor$. Phase 1 ($n<p\le2n$, $L=0$) is proven
in \cite{CBcert}. For $p\le n$ the following structure was discovered and
verified in this campaign.

\begin{finding}[verified; $272$ descents, $46$ band pairs, $0$ failures]
\label{find:dwork}
Let $a=\lfloor n/p\rfloor$. Then
\[
p^{5}\,\frac{p_n}{q_n}\;\equiv\;\frac{p_a}{q_a},
\qquad
p^{3}\,\frac{\wt p_n}{q_n}\;\equiv\;\frac{\wt p_a}{q_a}
\pmod{p^{k}},\qquad k\ge 2-\kappa\ \text{(typical }3-\kappa\text{)},
\]
for all tested $(n,p)$, $p\ge5$, $n\le45$. The digit constants are exact
ladder values --- in particular
\[
\frac{p_1}{q_1}=\frac{29}{28},\qquad
\frac{\wt p_1}{q_1}=\frac{101}{84},\qquad
\frac{p_2}{q_2}=\frac{24289}{23424},
\]
identified first as the unique low-height fit to the measured $p$-adic digits
and then confirmed exactly. The first correction digit is antisymmetric in
$r=n-ap$ about $(p-1)/2$ --- the derivative signature of a Dwork
$F(x)/F(x^p)$ congruence. Plain Lucas congruences
$q_{ap+b}\equiv q_aq_b\ (p)$ \emph{fail} at every tested prime: the operative
structure is a Frobenius descent of the ratio of solutions, not a digit
product. At exceptional primes dividing a digit-constant denominator (e.g.\
$7\mid28$), the ratio statement degrades but the combined integer invariant
$\ord_p(p_n)\ge\kappa-5L$ is preserved in every tested case.
\end{finding}

For the all-prime assembly it is useful to distinguish the observed ratio
congruence from the integer-lattice statement that it must induce. Put
$\kappa_p(m)=\ord_p\binom{2m}{m}$, and write $a=\lfloor n/p\rfloor$. The
required denominator-free form is the one-step inequality
\begin{equation}
 \tag{DWORK$_\Z$}\label{eq:dworkZ}
 \ord_p(p_n)-\kappa_p(n)
 \;\ge\;\ord_p(p_a)-\kappa_p(a)-5.
\end{equation}
Along a clean descent path, (DWORK) together with the free divisibility of
$q_n$ already gives the final bound directly.  The stronger one-step form
(DWORK$_\Z$) also records transport of any excess divisibility of $q_a$. At an
exceptional prime it asserts, in particular, that the excess divisibility of
$q_n$ compensates for the denominator of the relevant digit ratio. That
compensation is verified in the finite data but is part of the open theorem;
mere cross-multiplication of the displayed ratio congruence does not by itself
prove it.

\begin{proposition}[assembly; proven modulo (DWORK$_\Z$)]\label{prop:assembly}
Assume \eqref{eq:dworkZ} for every $p\ge5$ and every $n\ge p$. Then
\[
 \ord_p(p_n)\ge \kappa_p(n)-5\lfloor\log_p n\rfloor
 \qquad(p\ge5).
\]
Consequently Phase 2 of the central-binomial law holds. Combined with Phase
1 \cite{CBcert}, this yields $\binom{2n}{n}\mid6A_n$ and hence
$12\,d_n^5P_n\in\Z$ for all $n$.
\end{proposition}

\begin{proof}
Set $n_0=n$ and $n_{j+1}=\lfloor n_j/p\rfloor$. For
$L=\lfloor\log_p n\rfloor$ we have $n_L<p$. Applying
\eqref{eq:dworkZ} successively gives
\[
 \ord_p(p_n)-\kappa_p(n)
 \ge \ord_p(p_{n_L})-\kappa_p(n_L)-5L.
\]
Here $\kappa_p(n_L)=0$, and $p_{n_L}$ is $p$-integral because its known
denominator divides a power of $d_{n_L}$, with $p>n_L$. The right-hand side
is therefore at least $-5L$, proving the claim.
\end{proof}

Both the one-step form \eqref{eq:dworkZ} and the resulting inequality
$\ord_p(p_n)\ge\kappa-5L$ are verified with zero
failures for all $p\ge5$, $n\le45$ --- including the $206$ zero-slack tight
cases of the earlier ledger, which the descent now explains: the $d_n^5$
reserve pays exactly $5$ orders per digit level, and $\kappa$ carries the
Kummer carries.

\smallskip\noindent\textbf{Literature status and proof program.} A targeted
audit (with searches of the primary sources) indicates (DWORK) is \emph{not}
covered off-the-shelf: Delaygue--Rivoal--Roques \cite{DRR} treat
hypergeometric coefficient sequences and mirror-map quotients, not
harmonic-companion ratios; Malik--Straub \cite{MalikStraub} prove scalar
Lucas-type results. Derivative-corrected Lucas congruences for the classical
Ap\'ery sequence \cite{RYK} are a close precedent for the observed
antisymmetric correction, but still concern the leading sequence rather than
a coefficientwise ratio of independent solutions. Matrix Dwork theory
\cite{BeukersMatrix} supplies the correct higher-rank architecture, not an
off-the-shelf theorem for $p_n/q_n$. The
planned proof route: (1) restate (DWORK) denominator-free in the period
vector $(q_a,p_a,\wt p_a)$ and prove (DWORK$_\Z$), so exceptional primes
become explicit lattice bookkeeping;
(2) prove the $a=1$ band by splitting the very-well-poised summand into
$p$-blocks and expanding factorial and harmonic blocks through weight 5
(Wolstenholme-block cancellations); (3) treat $(p_n,\wt p_n)$ simultaneously,
exposing a triangular Frobenius matrix; (4) generalize in $a$. Phase 1 is the
proven terminal boundary condition for the iteration. It should not be called
the $a=0$ case of the displayed ratio descent: because of the prefactor $p^5$,
that formal congruence alone is too weak to recover Phase 1. To place the
weight-5 statement inside matrix Dwork theory
one still needs a period-matrix or diagonal realization of the complete
Zudilin triple. The period-level telescoping project for the distinct
weight-7 family (\S\ref{sec:rec}) may provide a useful model for that step, but
does not itself supply the required weight-5 realization.

\section{Outlook}\label{sec:outlook}

Three closures would upgrade the night's conditional results to theorems:
(1) a creative-telescoping certificate (or alternative proof) that the order-4
operator governs the weight-7 family at period level --- closing descent
membership, hence $s_3,\hat P_3$, hence hypothesis (ii) of the snap;
(2) a proof of the denominator-grid hypothesis (i), e.g.\ by the
Krattenthaler--Rivoal-type arithmetic for the weight-7 forms;
(3) a proof of (DWORK) together with its integer-lattice form
(DWORK$_\Z$), completing the corrected Brown--Zudilin law
$12d_n^5P_n\in\Z$. Independent of all three, the sharpest new open object is
the vector-valued Dwork descent itself: one congruence family, verified across
$272$ instances, whose archimedean limit is the classical convergence
$p_n/q_n\to\zt5$ and whose $p$-adic content is a self-similarity under
base-$p$ digit shift. One ratio, all completions, one mechanism.

\subsection*{Provenance}
Results obtained 17--18 July 2026 in a collaborative session: coordination,
reductions and proofs as cited; exact computation by fleet agents; independent
review by a second model (in particular the conditional-rigor form of
Proposition~\ref{prop:snap} and the literature audit of \S\ref{sec:dwork}).
Complete scripts and exact data in the repository (\texttt{worthiness/}):
recurrence and terms (\texttt{zeta7\_q\_recurrence.json},
\texttt{zeta7\_lc\_terms.txt}), snap and verification
(\texttt{ZETA7\_P3\_SNAP.md}, \texttt{ZETA7\_P3\_VERIFICATION.md},
\texttt{zeta7\_p3\_*.py}), sector measurement
(\texttt{zeta7\_sector\_measurement.py}), Phase-2 program
(\texttt{PHASE2\_BAND\_BLUEPRINT.md}, \texttt{PHASE2\_BAND\_V.md},
\texttt{lemma\_cb\_band\_v*.py}).

\begin{thebibliography}{9}
\bibitem{BZ} F.~Brown, W.~Zudilin, \emph{On cellular rational approximations
to $\zeta(5)$}, arXiv:2210.03391.
\bibitem{Zud02} W.~Zudilin, \emph{A third-order Ap\'ery-like recursion for
$\zeta(5)$}, Mat.\ Zametki \textbf{72} (2002); arXiv:math/0206178.
\bibitem{MOS} D.~McCarthy, R.~Osburn, A.~Straub, \emph{Sequences, modular
forms and cellular integrals}, arXiv:1705.05586.
\bibitem{CBcert} \emph{A Frobenius certificate for central-binomial
cancellation in Zudilin's symmetric $\zeta(5)$ construction}, companion note,
2026 (\texttt{worthiness/cb\_certificate.tex}).
\bibitem{Denoms} \emph{Sharp denominator corrections for Brown--Zudilin
cellular approximations}, companion paper, 2026
(\texttt{paper/denominators.tex}).
\bibitem{BarnesNotes} \emph{$\zeta(7)$ Barnes campaign notes}, 2026
(\texttt{worthiness/ZETA7\_BARNES.md}).
\bibitem{DRR} \'E.~Delaygue, T.~Rivoal, J.~Roques, \emph{On Dwork's $p$-adic
formal congruences theorem and hypergeometric mirror maps}, arXiv:1309.5902.
\bibitem{MalikStraub} A.~Malik, A.~Straub, \emph{Divisibility properties of
sporadic Ap\'ery-like numbers}, Res.\ Number Theory (2016).
\bibitem{RYK} E.~Rowland, R.~Yassawi, C.~Krattenthaler,
\emph{Lucas congruences for the Ap\'ery numbers modulo $p^2$}, Integers
\textbf{21} (2021), A20; arXiv:2005.04801.
\bibitem{BeukersMatrix} F.~Beukers, \emph{A matrix version of Dwork's
congruences}, arXiv:2005.01523.
\end{thebibliography}
\end{document}
